\subsubsection{Family Based Analysis}
\begin{enumerate}
    \item \textbf{Calculation of family Centroids}
          
          The function computes the centroid (mean expression vector) for each gene family
          in a gene expression dataset. It iterates over all families, selects the relevant
          genes for each family and calculates the mean vector for each family.
          There are two possible modes to use. To use all genes for calculation,
          select \texttt{all} mode. If you want to specify relevant genes for
          calculation, select \texttt{orthologs} mode. It is important that you provide
          a logical \texttt{ortholog\_set} array when using \texttt{orthologs} mode.
          Call the \texttt{tox\_group\_centroid} function with the following arguments:
          
          \begin{itemize}
            \item \texttt{expression\_vectors}: Expression vector matrix with size \texttt{n\_axes
                  $\times$ n\_genes}
            
            \item \texttt{gene\_to\_family}: Gene-to-family mapping array with the length
                  of \texttt{n\_genes} (1-based indexing)
                  
            \item \texttt{n\_families}: Total number of gene families
                  
            \item \texttt{mode}: Calculation mode ("all" or "orthologs")
                  
            \item (Optional) \texttt{ortholog\_set}: Logical array indicating if a gene
                  is part of the calculation subset (only used in \texttt{orthologs} mode)
          \end{itemize}
          
          \textbf{Return:}
          \begin{itemize}
            \item \texttt{centroid\_matrix}: Matrix of calculated family centroid vectors
                  with size \texttt{n\_axes $\times$ n\_families}
          \end{itemize}
          
          \begin{lstlisting}[language=R, caption={R Calculation of Family Centroids}]
tox_group_centroid(expression_vectors, gene_to_family, n_families, mode, 
    ortholog_set)
          \end{lstlisting}
          
          \begin{enumerate}
            \item \textbf{Calculate Mean Vector}
                  
                  Alternatively you can calculate a single mean vector for a given set of
                  vectors by calling the \texttt{mean\_vector} function directly with
                  the following arguments:
                  
                  \begin{itemize}
                    \item \texttt{expression\_vectors real(real64)}: Expression vector matrix
                          with size \texttt{n\_axes $\times$ n\_genes}
                          
                    \item \texttt{gene\_indices integer(int32)}: An array containing the
                          column indices of the selected genes in \texttt{expression\_vectors}
                  \end{itemize}
                  
                  \textbf{Return:}
                  \begin{itemize}
                    \item \texttt{centroid:} Calculated mean vector (centroid) with
                          length of \texttt{n\_axes}
                  \end{itemize}
                  \begin{lstlisting}[language=R, caption={R Mean Vector Calculation}]
tox_mean_vector(expression_vectors, gene_indices)
                  \end{lstlisting}
          \end{enumerate}
          
    \item \textbf{Compute Shift Vector Field}

    Calculates the shift vector field for each gene expression vector based on its family centroid. The shift vector represents the difference between the gene expression vector and its corresponding family centroid.
    
    \textbf{Mathematical Definition:}
    \[ \text{shift\_vector} = \text{expression\_vector} - \text{family\_centroid} \]
    
    \textbf{Arguments:}
    \begin{itemize}
        \item \texttt{expression\_vectors}: Matrix of size \texttt{(n\_axes\_genes × n\_vectors)} where each column is a gene expression vector
        \item \texttt{family\_centroids}: Matrix of size \texttt{(n\_axes\_centroids × n\_families)} where each column is a family centroid vector
        \item \texttt{gene\_to\_centroid}: Integer vector of length \texttt{n\_vectors} containing 1-based family indices mapping each gene to its centroid
    \end{itemize}
    \textbf{Returns:}
    \begin{itemize}
        \item \texttt{shift\_vectors}: Matrix of size \texttt{(2*n\_axes\_genes × n\_vectors)} containing centroids (rows 1:n\_axes\_genes) and shift vectors (rows (n\_axes\_genes+1):(2*n\_axes\_genes))
    \end{itemize}
    \begin{lstlisting}[language=R, caption={R Shift Vector Field Computation}]
shift_vectors <- tox_compute_shift_vector_field(expression_vectors, 
    family_centroids, gene_to_centroid)
    \end{lstlisting}          
    
    \item \textbf{Calculate Distance to Centroid}
          
          For each gene in the gene expression matrix, the function computes the
          Euclidean distance to its corresponding family centroid and returns a
          numeric vector of distances from each gene to its family centroid. Call the
          \texttt{tox\_distance\_to\_centroid} function with the following arguments:
          
          \begin{itemize}
            \item \texttt{genes}: Gene expression matrix with size of (\texttt{n\_genes
                  $\times$ d})
            
            \item \texttt{centroids}: Family centroid matrix with size of (\texttt{n\_families
                  $\times$ d})
            
            \item \texttt{gene\_to\_fam}: Gene-to-family mapping integer vector with
                  the length of \texttt{n\_genes} (0 = no family, \textgreater 0 =
                  family index)
                  
            \item \texttt{d}: Gene expression matrix dimension as integer
          \end{itemize}
          \begin{lstlisting}[language=R, caption={R Distance to Centroid Calculation}]
tox_distance_to_centroid(genes, centroids, gene_to_fam, d)
          \end{lstlisting}
          
          \begin{enumerate}
            \item \textbf{Calculate Euclidean Distance}
                  
                  Alternatively you can calculate a single Euclidean distance between two
                  vectors manually by calling \texttt{tox\_euclidean\_distance} function
                  with the following arguments:
                  
                  \begin{itemize}
                    \item \texttt{vec1:} First vector
                          
                    \item \texttt{vec2:} Second vector
                  \end{itemize}
                  
                  \textbf{Both vectors must have same length!}
                  
                  \begin{lstlisting}[language=R, caption={R Euclidean distance calculation}]
tox_euclidean_distance(vec1, vec2)
                  \end{lstlisting}
          \end{enumerate}
          
    \item \textbf{Detect Gene Outliers}
          
          The function first computes family scaling, then calculates relative
          distance indices (RDI) and identifies outliers based on the specified
          percentile threshold. It returns a list of components containing the outlier
          boolean values and intermediate results. Call the \texttt{tox\_detect\_outliers}
          function with the following arguments:
          
          \begin{itemize}
            \item \texttt{distances}: Numeric vector of distances for each gene to its
                  centroid
                  
            \item \texttt{gene\_to\_fam}: Gene-to-family mapping integer vector mapping
                  with the same length as the distances vector (0 = no family, \textgreater
                  0 = family index)
                  
            \item \texttt{n\_families}: Integer number of families
                  
            \item \textbf{(Optional)} \texttt{percentile}: Percentile threshold
                  value for outlier detection (default is 95.0)
          \end{itemize}
          
          \textbf{Return list:}
          \begin{itemize}
            \item \textbf{is\_outlier}: Logical vector indicating whether each gene in
                  an outlier
                  
            \item \textbf{loess\_x}: Numeric vector of x coordinates used as reference
                  points for LOESS smoothing (median distance)
                  
            \item \textbf{loess\_y}: Numeric vector of y coordinates used as reference
                  points for LOESS smoothing (standard deviation)
                  
            \item \textbf{loess\_n}: Integer vector with numbers of genes in each family
                  used to compute the corresponding median and standard deviation
          \end{itemize}
          
          \begin{lstlisting}[language=R, caption={R Gene Outlier Calculation}]
tox_detect_outliers(distances, gene_to_fam, n_families, percentile)
          \end{lstlisting}
          
          \textbf{Alternatively you can calculate each step of \texttt{tox\_detect\_outliers}
          manually with the following functions:}
          \begin{enumerate}
            \item \textbf{Compute Family Scaling}
                  
                  The function calculates a scaling factor for each gene family by computing
                  the median and standard deviation of gene-to-centroid distances within
                  each family, then applies LOESS smoothing to these values. Call the \texttt{tox\_compute\_family\_scaling}
                  function with the following arguments:
                  
                  \begin{itemize}
                    \item \texttt{distances}: Numeric vector of distances for each gene to
                          its centroid
                          
                    \item \texttt{gene\_to\_fam}: Gene-to-family mapping integer vector mapping
                          with the same length as the distances vector (0 = no family, \textgreater
                          0 = family index)
                          
                    \item \texttt{n\_families}: Integer number of families
                  \end{itemize}
                  
                  \textbf{Return dictionary:}
                  \begin{itemize}
                    \item \textbf{dscale}: Numeric vector of scaling factors per family
                          
                    \item \textbf{loess\_x}: Numeric vector of x coordinates used as reference
                          points for LOESS smoothing (median distance)
                          
                    \item \textbf{loess\_y}: Numeric vector of y coordinates used as reference
                          points for LOESS smoothing (standard deviation)
                          
                    \item \textbf{indices\_used}: Integer vector of indices of families that
                          have been included in the calculation
                  \end{itemize}
                  
                  \begin{lstlisting}[language=R, caption={R Compute Family Scaling}]
tox_compute_family_scaling(distances, gene_to_fam, n_families)
                  \end{lstlisting}
                  
            \item \textbf{Compute Family Scaling (Expert Version)}
                  
                  This function is the expert version of \texttt{tox\_detect\_outliers}
                  and requires pre-allocated work arrays for maximum performance and
                  control. It can be used when you need fine-grained control over memory
                  allocation or you are calling this function many times in a tight loop.
                  Call the \texttt{tox\_compute\_family\_scaling\_expert} function with the
                  following arguments:
                  
                  \begin{itemize}
                    \item \texttt{distances}: Numeric vector of distances for each gene to
                          its centroid
                          
                    \item \texttt{gene\_to\_fam}: Gene-to-family mapping integer vector mapping
                          with the same length as the distances vector (0 = no family, \textgreater
                          0 = family index)
                          
                    \item \texttt{n\_families}: Integer number of families
                          
                    \item \texttt{perm\_tmp}: Permutation integer vector for sorting with
                          the same length as distances
                          
                    \item \texttt{stack\_left\_tmp}: Stack integer vector for sorting with
                          the same length as distances
                          
                    \item \texttt{stack\_right\_tmp}: Stack integer vector for sorting with
                          the same length as distances
                          
                    \item \texttt{family\_distances}: Work integer vector for sorting with
                          the same length as distances
                  \end{itemize}
                  
                  \textbf{Return list:}
                  \begin{itemize}
                    \item \textbf{dscale}: Numeric vector of scaling factors per family
                          
                    \item \textbf{loess\_x}: Numeric vector of x coordinates used as reference
                          points for LOESS smoothing (median distance)
                          
                    \item \textbf{loess\_y}: Numeric vector of y coordinates used as reference
                          points for LOESS smoothing (standard deviation)
                          
                    \item \textbf{indices\_used}: Integer vector of indices of families that
                          have been included in the calculation
                          
                    \item \textbf{perm\_tmp}: Permutation vector for sorting
                          
                    \item \textbf{stack\_left\_tmp}: Stack vector for sorting
                          
                    \item \textbf{stack\_right\_tmp}: Stack vector for sorting
                          
                    \item \textbf{family\_distances}: Work vector for sorting
                  \end{itemize}
                  \begin{lstlisting}[language=R, caption={R Compute Family Scaling (Expert Version)}]
tox_compute_family_scaling_expert(distances, gene_to_fam, n_families, 
    perm_tmp, stack_left_tmp, stack_right_tmp, family_distances)
                  \end{lstlisting}
                  
            \item \textbf{Compute Relative Distance Index (RDI)}
                  
                  The function calculates the Relative Distance Index (RDI) for each gene
                  by dividing its Euclidean distance to the family centroid by the scaling
                  factor of its family. Call the \texttt{tox\_compute\_rdi} function with
                  the following arguments:
                  
                  \begin{itemize}
                    \item \texttt{distances}: Numeric vector of distances for each gene to
                          its centroid
                          
                    \item \texttt{gene\_to\_fam}: Gene-to-family mapping integer vector mapping
                          with the same length as the distances vector (0 = no family, \textgreater
                          0 = family index)
                          
                    \item \textbf{dscale}: Numeric vector of scaling factors per family
                  \end{itemize}
                  
                  \textbf{Return list:}
                  \begin{itemize}
                    \item \texttt{rdi}: Numeric vector containing the RDI value for each
                          gene
                          
                    \item \texttt{sorted\_rdi}: Numeric vector containing the RDI values
                          sorted in ascending order
                  \end{itemize}
                  
                  \begin{lstlisting}[language=R, caption={R Compute Relative Distance Index (RDI)}]
tox_compute_rdi(distances, gene_to_fam, dscale)
                  \end{lstlisting}
                  
            \item \textbf{Identify Outliers}
                  
                  The function identifies outliers based on the top percentile of RDI values.
                  Call the \texttt{tox\_identify\_outliers} function with the following arguments:
                  
                  \begin{itemize}
                    \item \texttt{rdi}: Numeric vector containing the RDI value for each
                          gene
                          
                    \item \textbf{(Optional)} \texttt{percentile}: Percentile threshold
                          number for outlier detection (default is 95.0)
                  \end{itemize}
                  
                  \textbf{Return list:}
                  \begin{itemize}
                    \item \texttt{is\_outlier}: Logical vector indicating whether each gene
                          is an outlier
                          
                    \item \texttt{threshold}: Actual threshold value used for outlier detection
                  \end{itemize}
                  \begin{lstlisting}[language=R, caption={R Identify Outliers}]
tox_identify_outliers(rdi, percentile)
                  \end{lstlisting}
          \end{enumerate}
          
    \item \textbf{Compute Normalized Tissue Versatility:}
          
          The function calculates a normalized tissue versatility score for selected
          gene expression vectors. It is based on the angle between each gene
          expression vector and the space diagonal. Versatility is normalized to [0,
          1], where 0 means uniform expression 1 means expression in only one axis.
          Call the \texttt{tox\_compute\_tissue\_versatility} function with the following
          parameters:
          
          \begin{itemize}
            \item \texttt{expression\_vectors}: Numeric matrix array with size \texttt{n\_axes
                  $\times$ n\_vectors}
            
            \item \texttt{vector\_selection}: Logical vector with length of \texttt{n\_vectors}
                  to specify which vectors should be included in the calculation
                  
            \item \texttt{axes\_selection}: Logical vector with length of \texttt{n\_axes}
                  to specify which axes should be included in the calculation
          \end{itemize}
          
          \textbf{Return list:}
          \begin{itemize}
            \item \texttt{tissue\_versatilities}: Numeric vector containing the calculated
                  tissue versatilities for selected vectors
                  
            \item \texttt{tissue\_angles\_deg}: Numeric vector containing the calculated
                  angles in degrees for selected vectors
                  
            \item \texttt{n\_selected\_vectors}: Integer number of vectors processed
                  
            \item \texttt{n\_selected\_axes}: Integer number of axes used in calculation
          \end{itemize}
          
          \begin{lstlisting}[language=R, caption={R Tissue Versatility Calculation}]
tox_calculate_tissue_versatility(expression_vectors, vector_selection, axis_selection)
          \end{lstlisting}
\end{enumerate}